\section{Discussion}
\label{sec:discussion}

\subsection{Role in UAV Safety}

\sys{} is an additional voice-facing safety interaction layer for small UAVs. It does not replace obstacle avoidance, geofencing, failsafe landing, manual override, or the flight controller. Those mechanisms still monitor the physical environment, enforce operating regions, handle vehicle faults, and preserve operator authority. \sys{} addresses a different interface problem: how nearby human speech can enter the UAV system as a constrained intent-state update rather than as an immediately executable command.

This distinction is important for ordinary movement speech. A movement intent state does not encode a complete navigation goal, direction, distance, speed, or permission to fly. It only records that the speaker may be requesting an ordinary UAV interaction, which must still be interpreted by a safety-state boundary and the vehicle policy. Similarly, an emergency intent state is a safety-critical signal that should receive priority handling, not proof that a physical hazard has been detected. Unknown is the containment state for unsupported or uncertain audio, keeping ambiguous speech from becoming an action request.

\subsection{Scope of the Baselines}

The transcript-first baseline in this paper is deliberately scoped. It uses Whisper \texttt{tiny.en} with a deterministic intent parser on fixed rotor-noisy one-second clips. This baseline represents a straightforward speech-to-text and parsing design under the same short-window acoustic condition used by the voice layer. It should not be read as a claim that all ASR systems fail around drones. Larger models, longer windows, different microphones, cloud processing, or task-specific ASR adaptation may change the result, but they also change the deployment and interaction assumptions.

The direct-mapping and no-unknown analyses have a different role. They are offline action-pressure simulations that ask what would happen if recognizer outputs were bound more directly to vehicle-facing actions, or if unsupported audio could no longer remain unknown. They are not flight tests, measured actuation results, or validated flight-safety evidence. Their purpose is to show why the intent-state mapping matters: errors and uncertainty should be contained before they reach a vehicle-facing decision.

The compact classifier baselines are also recognizer baselines, not complete safety-system baselines. They help compare acoustic model behavior and deployment feasibility, but they do not by themselves implement the intent-state mapping, unknown containment, or safety-state boundary that define \sys{} as a system.


\section{Conclusion}
\label{sec:conclusion}

\sys{} frames talk-to-the-drone interaction as a safety-state mediated way to admit human speech into UAV control. Instead of treating speech as a transcript, a keyword trigger, or a direct control input, the system converts short rotor-noisy audio windows into constrained intent-state updates and lets a safety-state boundary decide how those updates may influence the vehicle. This design keeps voice interaction natural and hands-free while preserving a conservative path between uncertain speech and UAV action.
